\section{引力波物理}
\subsection{谐和坐标条件的线性近似}

\subsubsection{场方程的自由度和规范选取}

考虑四维时空中的度规张量 $g_{\mu\nu}$ 是对称张量，一共有 $10$ 个独立分量。作为关于度规的场方程，爱因斯坦方程写为
\begin{align}
G_{\mu\nu} = 8\pi G\,T_{\mu\nu}\,,
\end{align}
其中 $G_{\mu\nu}$ 为爱因斯坦张量，$T_{\mu\nu}$ 为能量动量张量，$G$ 为万有引力常数。由于 $G_{\mu\nu}$ 也是一个对称二阶张量，表面上这给出了 $10$ 个标量方程。

然而实际上，Bianchi 恒等式
\begin{align}
\nabla_{\nu}G^{\mu\nu} = 0\,,
\end{align}
以及物质场的协变守恒
\begin{align}
\nabla_{\nu}T^{\mu\nu} = 0\,,
\end{align}
使得爱因斯坦方程的 $10$ 个分量中只有 $6$ 个是独立的：对 $G_{\mu\nu}=8\pi G T_{\mu\nu}$ 取协变散度后，上述两个恒等式保证其中有 $4$ 个分量是由其余分量自动满足的“兼容条件”。

\noindent
方程少了 $4$ 个真正独立的约束，正好与坐标（规范）自由度的 $4$ 个函数自由度相配合。广义相对论在任意坐标变换
\begin{align}
x^{\mu} \;\longrightarrow\; \widetilde{x}^{\mu} = \widetilde{x}^{\mu}(x)
\end{align}
下保持形式不变，这反映了理论的微分同胚不变性。将其写成无穷小形式，有
\begin{align}
x^{\mu} \;\longrightarrow\; \widetilde{x}^{\mu} = x^{\mu} + \xi^{\mu}(x)\,,
\end{align}
其中 $\xi^{\mu}(x)$ 是任意四个足够光滑的小函数，用来描述坐标变换的生成矢量场。因此，度规 $g_{\mu\nu}$ 的 $10$ 个分量中，有 $4$ 个自由方向实际上是坐标选择带来的 \emph{规范自由度}；只有在给定一定的坐标（规范）条件后，场方程才能具体决定度规的形式。

\paragraph{de Donder 规范与谐和坐标}

记 $g := \det(g_{\mu\nu})$，定义
\begin{align}
\Gamma^{\lambda} := g^{\mu\nu}\,\Gamma^{\lambda}{}_{\mu\nu}\,,
\end{align}
其中 $\Gamma^{\lambda}{}_{\mu\nu}$ 为 Levi--Civita 联络的 Christoffel 符号。可以证明
\begin{align}
\Gamma^{\lambda}
=
-\frac{1}{\sqrt{-g}}\,
\partial_{\rho}
\left(
  \sqrt{-g}\,g^{\rho\lambda}
\right)\,.
\end{align}
\begin{proof}
    
\begin{align}
\Gamma^{\lambda} &= g^{\mu\nu} \frac{1}{2} g^{\lambda\rho} \left( \partial_{\mu} g_{\rho\nu} + \partial_{\mu} g_{\rho\nu} - \partial_{\rho} g_{\mu\nu} \right) \\
&= g^{\lambda\rho} g^{\mu\nu} \partial_{\mu} g_{\rho\nu} - \frac{1}{2} g^{\mu\nu} g^{\lambda\rho} \partial_{\rho} g_{\mu\nu} \\
&= - g^{\rho\lambda} g_{\rho\nu} \partial_{\mu} g^{\mu\nu} - \frac{1}{2} g^{\rho\lambda} \frac{1}{g} \partial_{\rho} g \\
&= - \delta^{\rho}_{\nu} \partial_{\mu} g^{\mu\nu} - g^{\lambda\rho} \partial_{\rho} \ln \sqrt{-g} \\
&= - \frac{1}{\sqrt{-g}} \sqrt{-g} \partial_{\mu} g^{\mu\nu} - \frac{1}{\sqrt{-g}} g^{\rho\lambda} \partial_{\rho} \sqrt{-g} \\
&= - \frac{1}{\sqrt{-g}} \partial_{\mu} (\sqrt{-g} g^{\mu\nu})
\end{align}

\end{proof}


\begin{definition}[de Donder 规范]
    若坐标系满足
    \begin{align}
        \Gamma^{\lambda} := g^{\mu\nu}\,\Gamma^{\lambda}{}_{\mu\nu} = 0\,,
    \end{align}
    则称在该坐标系下采用了 \emph{de Donder 规范}（又称谐和坐标规范或引力的 Lorenz 规范）。
\end{definition}

对任意标量场 $\phi$，定义 d'Alembert 算子为
\begin{align}
\Box\phi
:=
\nabla_{\mu}\nabla^{\mu}\phi
=
\frac{1}{\sqrt{-g}}
\partial_{\rho}
\left(
  \sqrt{-g}\,
  g^{\rho\lambda}
  \partial_{\lambda}\phi
\right)\,.
\end{align}
利用上面关于 $\Gamma^{\lambda}$ 的结果，可以将其改写为
\begin{align}
\Box\phi
&=
\frac{1}{\sqrt{-g}}
\partial_{\rho}
\left(
  \sqrt{-g}\,
  g^{\rho\lambda}
\right)\partial_{\lambda}\phi
+
g^{\rho\lambda}\partial_{\rho}\partial_{\lambda}\phi
\\[0.3em]
&=
-\Gamma^{\lambda}\,\partial_{\lambda}\phi
+
g^{\rho\lambda}\partial_{\rho}\partial_{\lambda}\phi\,,
\end{align}
即
\begin{align}
\Box\phi
=
g^{\rho\lambda}\partial_{\rho}\partial_{\lambda}\phi
-
\Gamma^{\lambda}\,\partial_{\lambda}\phi\,.
\end{align}

\begin{derivationnote}[“坐标看作标量函数”的含义]
在流形上，一个坐标系可以看作一组映射
\begin{align}
x^{\mu} : M \longrightarrow \mathbb{R}\,,
\end{align}
其中 $\mu=0,1,2,3$。对每一个固定的 $\mu$，$x^{\mu}(p)$ 只是把点 $p$ 映到实数的一个函数，因此在张量意义上是一个标量场；其“上标 $\mu$”只是用来标记是哪一个坐标函数，而不是张量指标。这样，对于每个 $\mu$，都有
\begin{align}
\partial_{\lambda}x^{\mu} = \delta^{\mu}{}_{\lambda}\,,
\end{align}
而协变导数作用在标量上与普通偏导无异：
\begin{align}
\nabla_{\lambda}x^{\mu} = \partial_{\lambda}x^{\mu} = \delta^{\mu}{}_{\lambda}\,.
\end{align}
一般的标量场 $\phi(x)$ 不满足 $\partial_{\lambda}\phi = \delta^{\mu}{}_{\lambda}$ 这样的关系，只有坐标函数本身具有这一特殊性质。
\end{derivationnote}

现在可以自然地引入谐和坐标的概念。

\begin{definition}[谐和坐标]
若一组坐标函数 $x^{\mu}$（视作四个标量场）满足
\begin{align}
\Box x^{\mu} = 0\,,
\end{align}
则称 $x^{\mu}$ 为\emph{谐和坐标}。
\end{definition}

\begin{theorem}[谐和坐标与 de Donder 规范的等价性]
    取谐和坐标相当于挑选了de Donder规范。
\end{theorem}
\begin{proof}
    把 $\phi$ 取为坐标函数 $x^{\mu}$，利用
\begin{align}
\partial_{\lambda}x^{\mu} = \delta^{\mu}{}_{\lambda}\,,\qquad
\partial_{\rho}\partial_{\lambda}x^{\mu} = 0\,,
\end{align}
在
\begin{align}
\Box\phi
=
g^{\rho\lambda}\partial_{\rho}\partial_{\lambda}\phi
-
\Gamma^{\lambda}\,\partial_{\lambda}\phi
\end{align}
中代入，得到
\begin{align}
\Box x^{\mu}
&=
g^{\rho\lambda}\partial_{\rho}\partial_{\lambda}x^{\mu}
-
\Gamma^{\lambda}\,\partial_{\lambda}x^{\mu}
\\[0.3em]
&=
0
-
\Gamma^{\lambda}\,\delta^{\mu}{}_{\lambda}
=
-\,\Gamma^{\mu}\,.
\end{align}
    对任意坐标系，有恒等式
    \begin{align}
        \Box x^{\mu} = -\Gamma^{\mu}\,.
    \end{align}
    因此，谐和坐标条件
    \begin{align}
        \Box x^{\mu} = 0
    \end{align}
    与 de Donder 规范条件
    \begin{align}
        \Gamma^{\mu} = 0
    \end{align}
    完全等价。
\end{proof}
\subsubsection{弱场线性近似方程}

假设时空接近闵氏时空，
\begin{align}
g_{\mu\nu} = \eta_{\mu\nu} + h_{\mu\nu}\,,
\end{align}
其中 $\eta_{\mu\nu}$ 为闵氏度规，$h_{\mu\nu}$ 为小扰动（$|h_{\mu\nu}|\ll 1$）。指标的升降均使用 $\eta_{\mu\nu}$ 与 $\eta^{\mu\nu}$。在只保留一阶 $h_{\mu\nu}$ 的线性近似下，Christoffel 符号化为
\begin{align}
\Gamma^{\mu}{}_{\alpha\beta}
=
\frac{1}{2}\,
\eta^{\mu\nu}
\left(
  \partial_{\alpha}h_{\nu\beta}
  + \partial_{\beta}h_{\nu\alpha}
  - \partial_{\nu}h_{\alpha\beta}
\right)\,.
\end{align}
Ricci 张量 $R_{\mu\nu}$ 为
\begin{align}
R_{\mu\nu}
&=
\partial_{\lambda}\Gamma^{\lambda}{}_{\mu\nu}
-
\partial_{\nu}\Gamma^{\lambda}{}_{\mu\lambda}
\\[0.3em]
&=
-\frac{1}{2}
\Big(
  \partial_{\alpha}\partial^{\alpha}h_{\mu\nu}
  + \partial_{\mu}\partial_{\nu}h
  - \partial_{\alpha}\partial_{\mu}h^{\alpha}{}_{\nu}
  - \partial_{\alpha}\partial_{\nu}h^{\alpha}{}_{\mu}
\Big)\,,
\end{align}
其中
\begin{align}
h := \eta^{\mu\nu}h_{\mu\nu}
\end{align}
是扰动的迹，$\partial^{\alpha} := \eta^{\alpha\beta}\partial_{\beta}$。

为简化记号，引入\emph{迹反}操作:
\begin{definition}[迹反扰动张量]
度规扰动 $h_{\mu\nu}$ 的迹反定义为
\begin{align}
\overline{h}_{\mu\nu}
:=
h_{\mu\nu}
-
\frac{1}{2}\,
\eta_{\mu\nu}h\,.
\end{align}
\end{definition}

\begin{derivationnote}[四维中迹反操作的自反性]
在四维闵氏空间中，再次对 $\overline{S}_{\mu\nu}$ 做迹反，有
\begin{align}
\overline{\overline{S}}_{\mu\nu}
&=
\overline{S}_{\mu\nu}
-
\frac{1}{2}\eta_{\mu\nu}\,\overline{S}^{\alpha}{}_{\alpha}
\\[0.3em]
&=
\left(
  S_{\mu\nu}
  -
  \frac{1}{2}\eta_{\mu\nu}S^{\alpha}{}_{\alpha}
\right)
-
\frac{1}{2}\eta_{\mu\nu}
\left(
  S^{\alpha}{}_{\alpha}
  -
  2 S^{\alpha}{}_{\alpha}
\right)
\\[0.3em]
&=
S_{\mu\nu}\,,
\end{align}
这里用到了 $\eta^{\mu\nu}\eta_{\mu\nu}=4$。因此在四维中迹反操作是自反的：$\overline{\overline{S}}_{\mu\nu}=S_{\mu\nu}$。
\end{derivationnote}

\paragraph{谐和坐标条件的线性化形式}

在线性近似中，$\Gamma^{\mu}$ 写为
\begin{align}
\Gamma^{\mu}
=
g^{\alpha\beta}\Gamma^{\mu}{}_{\alpha\beta}
\simeq
\eta^{\alpha\beta}\Gamma^{\mu}{}_{\alpha\beta}
=
\partial_{\alpha}h^{\mu\alpha}
-
\frac{1}{2}\partial^{\mu}h\,,
\end{align}
其中 $h^{\mu\alpha}=\eta^{\mu\rho}\eta^{\alpha\sigma}h_{\rho\sigma}$。谐和坐标条件 $\Gamma^{\mu}=0$ 在线性阶给出
\begin{align}
\partial_{\alpha}h^{\mu\alpha}
-
\frac{1}{2}\partial^{\mu}h
=0\,.
\end{align}
用迹反扰动 $\overline{h}_{\mu\nu}$ 表示，注意
\begin{align}
\overline{h}^{\mu\alpha}
=
h^{\mu\alpha}
-
\frac{1}{2}\eta^{\mu\alpha}h\,,
\end{align}
可得谐和条件的等价形式
\begin{align}
\partial_{\alpha}\overline{h}^{\mu\alpha} = 0\,.
\end{align}
这就是说，在线性化的形式下，如果需要取谐和坐标，只需要满足线性扰动之逆反的四散度为0即可。
\paragraph{线性化爱因斯坦方程}

线性化的 Ricci 张量可以用迹反扰动 $\overline{h}_{\mu\nu}$ 重写，从而爱因斯坦张量线性部分 $G^{(1)}_{\mu\nu}$ 在谐和坐标下大幅简化。最终的线性化爱因斯坦方程可写为
\begin{align}
\overline{R}_{\mu\nu}
=
8\pi G\,T_{\mu\nu}\,,
\end{align}
整理后得到（原始形式仅作记录）
\begin{align}
\frac{1}{2}
\left(
  \partial_{\alpha}\partial^{\alpha}\overline{h}_{\mu\nu}
  +
  \eta_{\mu\nu}\partial_{\alpha}\partial_{\beta}\overline{h}^{\alpha\beta}
  -
  \partial_{\nu}\partial_{\alpha}\overline{h}_{\mu}{}^{\alpha}
  -
  \partial_{\mu}\partial_{\alpha}\overline{h}_{\nu}{}^{\alpha}
\right)
=
-8\pi G\,T_{\mu\nu}\,.
\end{align}
在谐和条件
\begin{align}
\partial_{\alpha}\overline{h}^{\mu\alpha}=0
\end{align}
下，含 $\partial_{\alpha}\overline{h}^{\mu\alpha}$ 的项全部消失，方程简化为
\begin{align}
\partial_{\alpha}\partial^{\alpha}\overline{h}_{\mu\nu}
=
-16\pi G\,T_{\mu\nu}\,.
\end{align}
记 d'Alembert 算子 $\Box := \partial_{\alpha}\partial^{\alpha}$，则
\begin{align}
\Box\,\overline{h}_{\mu\nu}
=
-16\pi G\,T_{\mu\nu}\,.
\end{align}

\begin{derivationnote}[与电磁学 Lorenz 规范的类比形式]
电磁势 $A_{\mu}$ 在 Lorenz 规范 $\partial_{\mu}A^{\mu}=0$ 下满足
\begin{align}
\Box A^{\mu}
=
-4\pi j^{\mu}\,,
\end{align}
其中 $j^{\mu}$ 为电流四矢量。线性化引力扰动 $\overline{h}_{\mu\nu}$ 在谐和坐标（也称 de Donder 规范、引力的 Lorenz 规范）下满足
\begin{align}
\Box\,\overline{h}_{\mu\nu}
=
-16\pi G\,T_{\mu\nu}\,,
\end{align}
两者在方程形式上非常相似。这种形式上的对应对于建立物理直觉有帮助：$\overline{h}_{\mu\nu}$ 在一定意义上扮演了“引力势”的角色，$T_{\mu\nu}$ 则扮演“源”的角色。
\end{derivationnote}

\paragraph{电动力学中的残余规范自由度}

\n{先在更熟悉的电磁场中看一眼“选完规范之后还剩下什么自由度”，再类比到引力。}

\begin{definition}[电磁规范变换与 Lorenz 规范]
设 $A_{\mu}$ 为电磁四势，$F_{\mu\nu}$ 为电磁场强张量，定义
\begin{align}
F_{\mu\nu}
=
\partial_{\mu}A_{\nu}
-
\partial_{\nu}A_{\mu}\,.
\end{align}
在规范变换
\begin{align}
A_{\mu} \;\longrightarrow\; A'_{\mu}
=
A_{\mu} + \partial_{\mu}\chi
\end{align}
下，$F_{\mu\nu}$ 不变，其中 $\chi$ 为任意标量函数，这体现了电磁场的规范不变性。

在四维闵氏时空中引入 Lorenz 规范条件
\begin{align}
\partial_{\mu}A^{\mu} = 0\,.
\end{align}
\end{definition}

在此规范下，四维 Maxwell 方程
\begin{align}
\partial_{\mu}F^{\mu\nu} = 4\pi j^{\nu}
\end{align}
可以写成对电磁势的波动方程
\begin{align}
\Box A^{\nu} = -4\pi j^{\nu}\,,
\end{align}
其中 d'Alembert 算子定义为
\begin{align}
\Box := \partial_{\mu}\partial^{\mu}\,.
\end{align}

\begin{derivationnote}[Lorenz 规范下的残余规范函数]
现在考虑在已经满足 Lorenz 规范 $\partial_{\mu}A^{\mu}=0$ 的基础上，再做一次规范变换
\begin{align}
A_{\mu} \;\longrightarrow\; A'_{\mu}
=
A_{\mu} + \partial_{\mu}\chi\,,
\end{align}
其中 $\chi$ 为尚未指定的标量函数。新的四势 $A'_{\mu}$ 是否依然满足 Lorenz 规范，取决于 $\chi$。

先计算
\begin{align}
\partial_{\mu}A'^{\mu}
&=
\partial_{\mu}\big(A^{\mu} + \partial^{\mu}\chi\big)
\\[0.3em]
&=
\partial_{\mu}A^{\mu}
+
\partial_{\mu}\partial^{\mu}\chi
\\[0.3em]
&=
\partial_{\mu}A^{\mu}
+
\Box\chi\,.
\end{align}
若原势 $A_{\mu}$ 已经满足 $\partial_{\mu}A^{\mu}=0$，为了使变换后的 $A'_{\mu}$ 仍满足 Lorenz 规范，即
\begin{align}
\partial_{\mu}A'^{\mu} = 0\,,
\end{align}
就必须有
\begin{align}
\Box\chi = 0\,.
\end{align}
因此，在 Lorenz 规范已经施加之后，并不是所有的 $\chi$ 都被禁止，而是仍然可以进行满足
\begin{align}
\Box\chi = 0
\end{align}
的规范变换
\begin{align}
A_{\mu} \;\longrightarrow\; A'_{\mu}
=
A_{\mu} + \partial_{\mu}\chi\,.
\end{align}
这些“剩下允许的”规范函数 $\chi$ 就对应于电磁场在 Lorenz 规范下的\emph{残余规范自由度}。
\end{derivationnote}

\n{直觉上：我们用 Lorenz 规范 $\partial_{\mu}A^{\mu}=0$ 已经削掉了一部分规范自由度，但仍可以在满足 $\Box\chi=0$ 的前提下稍微“调整” $A_{\mu}$ 而不破坏这一规范条件，这部分自由度就是“残余”的那一块。电磁学里是标量 $\chi$，在线性化引力里会变成四分量的 $\xi^{\mu}$。}

\paragraph{无穷小坐标变换与残余规范自由度}

\n{即便已经选取了谐和坐标条件，仍然存在一些进一步的坐标变换不会破坏这一条件，它们对应的就是“残余规范自由度”。}

考虑无穷小坐标变换
\begin{align}
x^{\mu}
\;\longrightarrow\;
\widetilde{x}^{\mu}
=
x^{\mu}
+
\xi^{\mu}(x)\,,
\end{align}
其中 $\xi^{\mu}(x)$ 是小量。视此为\emph{被动}坐标变换，度规在两组坐标中的分量满足
\begin{align}
\widetilde{g}_{\mu\nu}(\widetilde{x})
=
\frac{\partial x^{\rho}}{\partial \widetilde{x}^{\mu}}
\frac{\partial x^{\sigma}}{\partial \widetilde{x}^{\nu}}
g_{\rho\sigma}(x)\,.
\end{align}
在弱场近似下，取
\begin{align}
g_{\mu\nu} = \eta_{\mu\nu} + h_{\mu\nu}\,,\qquad
\widetilde{g}_{\mu\nu} = \eta_{\mu\nu} + \widetilde{h}_{\mu\nu}\,,
\end{align}
并只保留 $\xi^{\mu}$ 和 $h_{\mu\nu}$ 的一阶项，可得到扰动的规范变换
\begin{align}
\widetilde{h}_{\mu\nu}
=
h_{\mu\nu}
-
\partial_{\mu}\xi_{\nu}
-
\partial_{\nu}\xi_{\mu}\,,
\end{align}
其中 $\xi_{\mu} := \eta_{\mu\nu}\xi^{\nu}$。这是在闵氏背景下李导数的线性化结果。

迹反扰动在该变换下的变化为
\begin{align}
\delta\overline{h}_{\mu\nu}
&:=
\widetilde{\overline{h}}_{\mu\nu}-\overline{h}_{\mu\nu}
\\[0.3em]
&=
-\partial_{\mu}\xi_{\nu}
-\partial_{\nu}\xi_{\mu}
+
\eta_{\mu\nu}\partial_{\alpha}\xi^{\alpha}\,.
\end{align}
取散度：
\begin{align}
\partial^{\nu}\delta\overline{h}_{\mu\nu}
&=
-\partial^{\nu}\partial_{\mu}\xi_{\nu}
-\partial^{\nu}\partial_{\nu}\xi_{\mu}
+
\partial_{\mu}\partial_{\alpha}\xi^{\alpha}
\\[0.3em]
&=
-\partial_{\mu}\partial_{\nu}\xi^{\nu}
-\Box\xi_{\mu}
+
\partial_{\mu}\partial_{\nu}\xi^{\nu}
\\[0.3em]
&=
-\Box\xi_{\mu}\,.
\end{align}

如果原本的扰动满足谐和条件
\begin{align}
\partial^{\nu}\overline{h}_{\mu\nu} = 0\,,
\end{align}
则经过变换后的扰动满足
\begin{align}
\partial^{\nu}\widetilde{\overline{h}}_{\mu\nu}
&=
\partial^{\nu}\overline{h}_{\mu\nu}
+
\partial^{\nu}\delta\overline{h}_{\mu\nu}
=
-\Box\xi_{\mu}\,.
\end{align}
要使新的扰动仍处于谐和规范中，必须有
\begin{align}
\partial^{\nu}\widetilde{\overline{h}}_{\mu\nu} = 0
\quad\Longrightarrow\quad
\Box\xi_{\mu} = 0\,.
\end{align}

\begin{definition}[谐和规范下的残余规范变换]
在谐和坐标条件 $\partial_{\nu}\overline{h}^{\mu\nu}=0$ 已经被施加的前提下，仍允许进行满足
\begin{align}
\Box\xi^{\mu} = 0
\end{align}
的无穷小坐标变换
\begin{align}
x^{\mu}\;\longrightarrow\;\widetilde{x}^{\mu}
=
x^{\mu} + \xi^{\mu}(x)\,,
\end{align}
并保持谐和条件不变。这些变换对应的自由度称为\emph{谐和规范下的残余规范自由度}。
\end{definition}

\n{这一结论与电磁学中 Lorenz 规范下的残余规范变换完全平行：在 $\partial_{\mu}A^{\mu}=0$ 条件下，$A_{\mu}\to A_{\mu}+\partial_{\mu}\chi$ 仍保持规范，只要 $\chi$ 满足 $\Box\chi=0$。在引力的线性理论中，$\xi^{\mu}$ 扮演了电磁学中标量规范函数的角色。}

综上，在弱场线性近似下：
\begin{align}
g_{\mu\nu}
&=
\eta_{\mu\nu} + h_{\mu\nu}\,,
\\[0.3em]
\overline{h}_{\mu\nu}
&=
h_{\mu\nu}
-
\frac{1}{2}\eta_{\mu\nu}h\,,
\\[0.3em]
\partial_{\nu}\overline{h}^{\mu\nu} &= 0\quad\text{（谐和坐标条件）}\,,
\\[0.3em]
\Box\overline{h}_{\mu\nu}
&=
-16\pi G\,T_{\mu\nu}\quad\text{（线性化爱因斯坦方程）}\,,
\\[0.3em]
\widetilde{h}_{\mu\nu}
&=
h_{\mu\nu}
-
\partial_{\mu}\xi_{\nu}
-
\partial_{\nu}\xi_{\mu}\,,
\\[0.3em]
\Box\xi^{\mu}
&=
0\quad\text{（保持谐和规范的残余自由度）}\,.
\end{align}
这组关系构成了“谱和谐和坐标条件线性近似”中最常用的一套工作公式。
